\documentclass[11pt]{article}
\usepackage[margin=1in]{geometry}
\usepackage{amsmath,amssymb,amsthm,mathtools,array,booktabs,longtable}
\usepackage{microtype,xcolor}
\usepackage[colorlinks=true,linkcolor=blue!55!black,citecolor=blue!55!black,urlcolor=blue!55!black]{hyperref}
\hypersetup{pdftitle={Gluing relative corners to decorated displays},pdfsubject={Relative pseudodistributive laws from independently supplied content and occurrence lifts}}
\newtheorem{theorem}{Theorem}[section]
\newtheorem{proposition}[theorem]{Proposition}
\newtheorem{lemma}[theorem]{Lemma}
\newtheorem{corollary}[theorem]{Corollary}
\theoremstyle{definition}\newtheorem{definition}[theorem]{Definition}
\newcommand{\B}{\mathcal B}\newcommand{\J}{\mathcal J}
\newcommand{\D}{\mathbb D}\newcommand{\N}{\mathbb N}
\newcommand{\K}{\mathcal K}\newcommand{\Lcat}{\mathcal L}
\newcommand{\V}{\mathcal V}\newcommand{\Set}{\mathsf{Set}}\newcommand{\Cat}{\mathsf{Cat}}
\newcommand{\id}{\mathrm{id}}\newcommand{\op}{\mathrm{op}}\newcommand{\dom}{\mathrm{dom}}
\newcommand{\cod}{\mathrm{cod}}\newcommand{\iso}{\xrightarrow{\ \sim\ }}
\newcommand{\ea}{\mathsf e}\newcommand{\un}{\mathbf1}\newcommand{\Tt}{\overline T}
\newcommand{\Ssp}{\mathsf S_{\J}}\newcommand{\Disc}{\mathrm{disc}}
\newcommand{\st}{\mathbin{\star}}\newcommand{\bx}{\mathop{\boxtimes}}
\newcommand{\pure}{\mathsf{p}}\newcommand{\grade}{\mathsf G}
\newcommand{\pb}{\mathrm{pb}}\newcommand{\A}{\mathcal A}
\title{Gluing relative corners to decorated displays\\
\large A law of relative modules from an independently lifted occurrence monad}
\author{Contextual Composition project technical note}\date{20 November 2026}
\begin{document}\maketitle
\noindent This technical note was developed with substantial AI assistance.
The proof is written; the companion script checks finite examples and is not
a formal certificate of the general theorem.
\begin{abstract}
A content arrow carries a Kleisli arity. Pulling a decorated display
back against that arity and pulling the actual content arrow back at
its codomain are different operations. A cartesian comparison of
their \emph{underlying squares} is the interface that allows the
two operations to be glued without identifying their native arrows.
The result is a pseudodistributive law over the entire pure category.
The four comparisons and their ten pastings are derived from a
decoration-only lifted monad, cartesian uniqueness on content, and
a pure pullback condition. Finite bispan evaluation with directed
additive transfer gives one lifted monad; a two-stage presheaf
example changes the polynomial, merges witnesses and forgets a
colour across a noninvertible restriction. No unmarked recognition
or free universal property for the assembly is asserted.
\end{abstract}

\paragraph{What is proved.}
Theorem~\ref{thm:glue} is a sufficient construction on an ordinary base
with category-valued (not groupoid-valued) modules. It does not posit an
intertwiner on mixed words. Its two inputs can be built and mapped
independently: a relative cartesian calculus on the content double
category and an arity action on decorated displays. Their common datum is
a cartesian monad on the base, not a common category of grade morphisms.
The chosen content lifts must map to the specified \emph{display
pullbacks}; the arity need not detect equality of content maps.
The characteristic nontrivial equation on decorations is
\eqref{eq:middle}; it is proved before there is a content module.
Sections~\ref{sec:set} and~\ref{sec:diagram} construct, rather than
assume, the decoration input in two models. Universality in the sense
of an initial object among extensions is a different claim
(Section~\ref{sec:boundary}).

\section{Two independent inputs}\label{sec:inputs}
Use nested universes in which the following ordinary categories and
category-valued pseudocoends exist. Let $\B$ have pullbacks, let $\J$
contain the isomorphisms and be stable under composition, pullback
and isomorphism of arrows. Let $(T,\eta,\mu)$ be a pullback-preserving
monad with cartesian unit and multiplication and $T\J\subseteq\J$.
Composition of arrows is written $vu=v\circ u$; the order of a word
of horizontal arrows or modules is \emph{execution order}. A Kleisli
arrow $a:A\to TB$ followed by $b:B\to TC$ has composite
$b\bullet a=\mu_C(Tb)a$. No restriction is placed on pure maps in $\B$.

\paragraph{Decorated displays.}
Take a pseudo-double category $\D$ with vertical category $\B$ and a
forgetful pseudo-double functor to the double category of commuting
squares in $\B$, identity on the vertical category. A horizontal arrow
$P:E\to B$ lies over a named display $r:E\to B$ in $\J$. A map
$P\to P'$ is a square over $(h:E\to E',u:B\to B')$, $r'h=ur$;
many different maps can lie over that same square. Horizontal composition
is written $P\st S:F\to B$ for $S:F\to E$. It lies over $rs$.
The horizontal unit $\un_B$ lies over $\id_B$. All associators and
unitors have identity boundary; as usual one may transport a named
carrier by an isomorphism. For definiteness use literal carrier
composition and chosen pullbacks; none of the results depend on this
choice.
Require $\cod:\D_1\to\B$ to be a fibration for \emph{all}
base maps, its chosen cartesian squares to be underlying pullbacks,
and its horizontal units and horizontal pastes of cartesian squares
to be cartesian. Denote a restriction by $u^*P$. This is a condition on
the available \emph{context} changes of base, not on relative content
lifts. The category $\K_B$ has objects the $P:E\to B$ and arrows the
squares over $u=\id_B$; $\ea_B(P)=E$ also acts on arrows by $h$.
Its associated module is
\[
 M(A,B)=(\ea_A\downarrow B),\qquad
 (P,g): A\xleftarrow{r} E\xrightarrow{g} B.
\]
A native map over $(k,v)$ is a decorated square over $(h,k)$ with
$vg=g'h$. The unit is $(\un_A,\id_A)$, composition of $(P,g)$
then $(S,f)$ uses $P\st g^*S$ and the induced forward tail.
These give a relative pseudopromonad: in a triple composite the
two reindexed last displays are the two cartesian lifts of the same
display along the same composite, after the canonical pullback
identifications. The horizontal associator of $\D$ then identifies the
two accumulations. The pentagon and triangles follow from cartesian
uniqueness on reindexings and the horizontal pentagon and triangles
of $\D$ on the resulting common diagram. This argument is natural in
all decorated squares by cartesian factorization; it does not assert
that the undecorated carrier detects their equality.

\begin{definition}[An independently lifted occurrence monad]\label{def:lift}
On this single double category take the following data.
\begin{enumerate}
\item A pseudo-double endofunctor $\Tt$ over $T$, whose object and vertical
parts are $T$, whose value on $P$ lies over $Tr:TE\to TB$, and which
preserves cod-cartesian squares. Its invertible horizontal comparisons
are denoted
$\Tt(P\st S)\cong\Tt P\st\Tt S$ and
$\Tt\un_B\cong\un_{TB}$ (colax orientation). They are natural on all
squares and obey the horizontal functor pentagon and the two unit
equations.
\item Vertical double transformations $\bar\eta:1\Rightarrow\Tt$,
$\bar\mu:\Tt^2\Rightarrow\Tt$ over $\eta,\mu$. Their squares
on $P$ lie over $(\eta_E,\eta_B)$ and $(\mu_E,\mu_B)$ and are
cod-cartesian in $\D_1$. A vertical double transformation is a
natural transformation on the vertical category and on the horizontal
arrow category, commuting with domain/codomain, the horizontal unit
and horizontal composition with the functor comparisons inserted.
Require the ordinary monad equalities of these vertical
transformations: $\bar\mu\,\Tt\bar\mu=\bar\mu\,\bar\mu_{\Tt}$ and
$\bar\mu\,\Tt\bar\eta=1=\bar\mu\,\bar\eta_{\Tt}$.
\end{enumerate}
These are equalities of maps of \emph{decorated displays}.
In particular the multiplication equation on two composable displays
is not an assertion about a content multiplication or grade.
This is a sufficient indexed replacement for addressed polynomial
evaluation when no finite-set polynomial presentation is chosen.
\end{definition}

\paragraph{Content and bare corners.}
Independently take a relative pseudopromonad
$N:\B^{\op}\times\B\to\Cat$ with its own multiplication and unit.
Equivalently for the present construction take its realized bicategory,
marked by the identity-on-objects pure pseudofunctor
$\pure:\B\to N$; native cells are the original hom-category arrows.
Pure pre- and postcomposition are written $k^*,u_*$, respectively,
and identified with the given actions by their unit comparisons.
Write $\N$ for the double category with horizontal arrows
$n\in N(A,B)$ and squares over $(k:A\to A',u:B\to B')$ the native
cells $u_*n\to k^*n'$. Vertical and horizontal compositions use the
given actions, horizontal composition in $N$, and the pure constraints.
In particular $n\to u_*n$ with left side $\id_A$ is opcartesian for
right-edge transport: factor a square with right side $vu$ by using
$(vu)_*\cong v_*u_*$. This requires no representability of $N$.
For details of this realization see \cite[Secs.~2--3]{U}.

Take an arity assignment $a_n:A\to TB$, preserving pure arrows,
horizontal multiplication and all squares:
\[
 a_{\pure(u)}=\eta u,\quad a_{nm}=a_m\bullet a_n,\qquad
 Tu\,a_n=a_{n'} k\quad\text{on every square } n\to n'.
 \tag{A}\label{eq:arity}
\]
Here $nm$ denotes $n$ followed by $m$. Equivalently this is a
pure-compatible multiplicative map to the locally discrete Kleisli
category. Structural isomorphisms therefore project to equalities.

\begin{definition}[Relative exactness for content]\label{def:corner}
For $n:A\rightsquigarrow B$ and $r:E\to B$ in $\J$ choose a square
\[
 C_r(n)=
 \left[
 \begin{array}{ccc}
 A_r&\xrightarrow{\ n_r\ }&E\\[-2pt]
 {\scriptstyle\ell}\downarrow&&\downarrow{\scriptstyle r}\\[-2pt]
 A&\xrightarrow{\ n\ }&B
 \end{array}\right].
\]
Require the following independent conditions on $\N\to\mathrm{Kl}(T)$.
\begin{enumerate}
\item $C_r(n)$ is cartesian for $\cod:\N_1\to\B$. Given any square
$s:z\to n$ with right side $rv$ there is exactly one square
$t:z\to n_r$ with right side $v$ and $C_r(n)t=s$.
Neither $v$ nor either left edge is required to lie in $\J$.
\item Its arity square is a pullback:
$(\ell,a_{n_r}):A_r\cong A\times_{TB} TE$ over $(a_n,Tr)$.
Thus $\ell\in\J$. Horizontal composites of these cartesian squares
are cartesian whenever matched: first lift the last horizontal
arrow and use its left edge to lift its predecessor.
Horizontal identities on a display are cartesian.
Preservation here concerns the \emph{actual} horizontal arrows of $N$.
\item For every pullback square
$E_u\xrightarrow{v}E,\ E_u\xrightarrow{r_u}B,\ E\xrightarrow r B',
\ B\xrightarrow u B'$ with $r\in\J$, the square of pure horizontal
arrows $\pure(v)\to\pure(u)$ over $(r_u,r)$ is cod-cartesian.
Equivalently, in the presence of the preceding clauses the canonical
base-change factors in \eqref{eq:bc} are invertible for all $u$.
\end{enumerate}
The square in the last clause is the pure functor's comparison of
$r v=u r_u$; it is not an arbitrarily selected isomorphic square.
This pure-exactness test involves no decoration. It implies the
horizontal-identity case by taking $u$ an identity.
\end{definition}

\section{Local calculus and the gluing theorem}\label{sec:glue}
For an arbitrary Kleisli arrow $a:A\to TB$ and $P:E\to B$ define
\[
 L_a P=a^*\Tt P,\qquad
 E_a=A\times_{TB}TE,\qquad a_P:E_a\to TE.
 \label{eq:La}\tag{1}
\]
The last map is again a Kleisli arrow, now to $E$.
For the horizontal unit write $I_a:L_a\un_B\cong\un_A$.
For composable Kleisli arrows $a,b$ and a display $P$ over the target
of $b$ write $K_{a,b}(P):L_a L_b P\cong L_{b\bullet a}P$.
For a pure $u$ write $U_u(P):L_{\eta u}P\cong u^*P$.

\begin{lemma}[Local indexed equations]\label{lem:local}
$L$ is functorial on decorated maps in each fibre and on changes of
Kleisli arrows by arbitrary commuting squares. In addition to $I,K,U$
there is an invertible comparison
\[
 F_a(P,S):L_a(P\st S)\ \iso\ (L_aP)\st L_{a_P}S.
 \tag{2}\label{eq:F}
\]
All are natural in maps of their typed input diagrams, including
noninvertible decorated maps. The $K,U$ obey the contravariant
pseudofunctor equations on Kleisli composites and units.
The $F,I$ obey horizontal associativity and units. The two structures
are compatible; in particular
\begin{align}
 F_{b\bullet a}(P,S)K_{a,b}(P\st S)
  ={}&
  \bigl(K_{a,b}(P)\st K_{a_{L_bP},b_P}(S)\bigr)
  \nonumber\\[-2pt]&{}\cdot
  F_a(L_bP,L_{b_P}S)\, L_a(F_b(P,S)).
 \tag{3}\label{eq:middle}
\end{align}
On the right, transport along the source isomorphism of $K_{a,b}(P)$
is understood in its second slot. Its defining arity equation is
$b_P\bullet a_{L_bP}=(b\bullet a)_P\,\dom K_{a,b}(P)$.
There are the corresponding horizontal-unit and pure-unit equations.
\end{lemma}
\begin{proof}
For $P\to P'$ over $(h,u)$ and $a,a'$ satisfying $Tu\,a=a'k$,
apply $\Tt$ and precompose the square $a^*\Tt P\to\Tt P$;
factor into the cartesian square for $a'^*\Tt P'$. Its left edge
has coordinates $(k,Th)$. Cartesian uniqueness gives functoriality
on actual maps, not just edges. For $K$ the chain is
\[
 a^*\Tt(b^*\Tt P)\cong a^*(Tb)^*\Tt^2P
 \cong (Tb\,a)^*(\mu^*\Tt P)\cong L_{b\bullet a}P .
\]
The first step uses cartesian preservation; the second uses the
cartesian square $\bar\mu_P$. On top projections it is exactly the
Kleisli equation in the statement. Use $\bar\eta_P$ for $U$.
For $F$, first use the compositor for $\Tt$; cartesian restriction
of $\Tt P\st\Tt S$ along $a$ is the horizontal composite of the
cartesian restriction of $\Tt P$ along $a$ and of $\Tt S$ along
$a_P$. This uses horizontal cartesian stability in $\D$. The horizontal
unit comparison gives $I$ in the same way.

Here is the equality test for these maps. Restore the arrows of each
restriction to its original display. In $K$-associativity, the two
total decorated arrows out of the triple lift then end by
$\bar\mu_P\Tt\bar\mu_P$ and $\bar\mu_P\bar\mu_{\Tt P}$;
naturality moves any intervening cartesian projections to their common
source. Equality follows by the monad multiplication equation.
The two $K$-unit equations use its two unit equations.
For $F$-associativity the restored arrows are the two horizontal
functoriality composites for $\Tt$ on three displays, and analogously
for the two units. For \eqref{eq:middle}, after restoration the
square is the following one, before restricting along $(Tb)a$:
\[
\begin{array}{ccc}
 \Tt^2(P\st S)&\longrightarrow&\Tt^2P\st\Tt^2S\\
 \bar\mu_{P\st S}\downarrow&&\downarrow\bar\mu_P\st\bar\mu_S\\
 \Tt(P\st S)&\longrightarrow&\Tt P\st\Tt S .
\end{array}
\]
The top comparison is the composite double-functor comparison.
The square commutes by the horizontal-composition axiom for the vertical
double transformation $\bar\mu$. The horizontal-unit version uses
its unit axiom; the versions with a pure occurrence unit use the
horizontal-composition and unit axioms for $\bar\eta$.
Cartesianness of the final restrictions makes equality after restoration
imply equality of the factors. This cancellation is in $\D_1$;
no fidelity of a coefficient-forgetting functor is used.
\end{proof}

\paragraph{The target in which equality is tested.}
For a module $X$ put $\grade_X(A,C)$ equal to the category with objects
$(P,x)$, $P\in\K_A$, $x\in X(\ea_AP,C)$, and arrows
\[
 (\alpha:P\to P',\quad t:x\to h^*x'),\qquad h=\ea_A\alpha.
 \tag{4}\label{eq:G}
\]
For moving endpoints $(k,v)$ use a decorated square over $(h,k)$
and $t:v_*x\to h^*x'$. Composition is the composition of the actual
arrows with the pseudofunctor comparisons. There is a module equivalence
\[
 M\odot X\simeq \grade_X,\qquad
 [Z;(P,g:E\to Z),x]\longmapsto(P,g^*x).
 \tag{5}\label{eq:normal}
\]
Indeed the relative tensor is the category-valued pseudocoend over
\emph{all} $Z\in\B$: generators are pairs of native arrows and invertible
balances along every pure arrow, with naturality, identity and composition
relations. A quasi-inverse chooses $Z=E$ and tail $\id_E$; on
\eqref{eq:G} first apply $t$, then balance along $h$, then apply $\alpha$.
Balance along $g$ provides the comparison to any representative, and the
balance relations give inverse equivalence and endpoint naturality.
In particular equality in \eqref{eq:G} is tested separately on
$\alpha$ and $t$, not on their underlying $\B$-maps.

\begin{proposition}[The local pullback being glued]\label{prop:matched}
Write $0$ as a superscript for undecorated $\J$-displays (with all
commuting squares), and $K_T(A,C)=\Disc\B(A,TC)$. At each $A,C$
the following is a pullback of categories in the common named-carrier
presentation, and a pseudopullback up to equivalence:
\[
\begin{array}{ccc}
 \grade_N^{\D}&\longrightarrow&\grade_N^{0}\\
 \downarrow&&\downarrow\\
 \grade_{K_T}^{\D}&\longrightarrow&\grade_{K_T}^{0}.
\end{array}
 \tag{5a}\label{eq:matched}
\]
Horizontal maps forget the decoration, vertical maps take arity.
\end{proposition}
\begin{proof}
A pair of objects of the top-right and bottom-left categories agrees
at bottom right precisely if it gives the same display $r:E\to A$
and the same arrow $a_n:E\to TC$. Such a pair is uniquely $(P,n)$.
A pair of arrows matches precisely if the left carrier map is the
same $h$; both the native decorated arrow over $h$ and the native
cell $n\to h^*n'$ then remain as \eqref{eq:G}. These bijections
respect compositions, including the restriction comparison used by
both copies of $\grade_N$. The right vertical functor is an
isofibration: an isomorphism of the underlying displays and arities
$h:(r,a_n)\cong(r',a')$ lifts by transporting the source of $n$ to
$(h^{-1})^*n$, which has exactly arity $a'$ and has the canonical
invertible comparison over $h$. Consequently any isomorphism used for
matching in the iso-comma category can be rectified to equality;
on already matched objects its morphism condition is the strict
pullback equation. This proves the pseudopullback assertion.
\end{proof}
This local mapping property explains why matching only the carrier of an
\emph{object} would be insufficient. To glue a family of comparison
maps both projected arrows have to agree in the lower-right category.
It does not supply any of the comparisons whose construction follows.

\begin{theorem}[Relative corners--decoration gluing]\label{thm:glue}
Under the hypotheses of Section~\ref{sec:inputs} there is a
pseudodistributive law of relative pseudopromonads
\[
                   \tau:N\odot M\longrightarrow M\odot N.
\]
On $[B;n:A\rightsquigarrow B,(P,g:E\to C)]$, $P$ over $r:E\to B$,
its normalized value is
\[
             (L_{a_n}P,\quad g_* n_r)\in\grade_N(A,C).
 \tag{6}\label{eq:cut}
\]
The notation transports the source of $n_r$ by the unique arity
pullback isomorphism in Definition~\ref{def:corner}(2). Its pure
balance is the cartesian--opcartesian mate, its structural isomorphisms
are those in \eqref{eq:four}, and it is natural on the category of
independent input maps in Section~\ref{sec:maps}. Consequently
$\grade_N$ carries the composite relative pseudopromonad, with
unit $j_M\odot j_N$ and multiplication
$(m_M\odot m_N)(M\odot\tau\odot N)$ after normalization.
\end{theorem}
\begin{proof}
\emph{Squares.} Take any square $s:n\to n'$ in $\N$ over $(k,u)$
and a decorated square $P\to P'$ over $(h,u)$. Factor in the
\emph{total} category of horizontal content arrows:
\[
 C_{r'}(n')\,\widetilde s=s\,C_r(n),\qquad
       \text{right side of }\widetilde s=h.
 \tag{7}\label{eq:square}
\]
Existence and uniqueness use relative cartesianness with unrestricted
test $h$. Write $H$ for its left edge. By the arity square for
$\widetilde s$ and the equation with $C_{r'}(n')$, $H$ has the two
coordinates $(k\ell,Th\,a_{n_r})$ in the target pullback.
The same two coordinates are the boundary of the \emph{decorated}
square obtained by applying Lemma~\ref{lem:local} to $P\to P'$.
Thus these two actual squares fit. For tails $g,g'$ and target map $v$
with $vg=g'h$, append pure squares, or equivalently write the bent
cell as
\[
 v_*g_* n_r \ \cong\ g'_*h_* n_r
 \longrightarrow g'_*H^*n'_{r'}\ \cong\ H^*g'_*n'_{r'} .
\]
This is generally not invertible. Vertical composition and identity
in \eqref{eq:square} follow by postcomposing with $C_{r'}(n')$,
using its cartesian uniqueness; the decorated maps compose by
Lemma~\ref{lem:local}. In particular all fixed-middle native product
arrows act. For a change of left endpoint $k$ the corresponding
comparison to restriction is invertible: lift the pure arrow
$\pure(k)$ through the left display of $C_r(n)$ by its pure pullback
square, and paste horizontally with $C_r(n)$. The result is cartesian
by horizontal stability and has bottom $k^*n$, up to the pure action
identification. Its decorated comparison is ordinary restriction of
\eqref{eq:La}. Right endpoint changes append pure arrows by the
pseudoaction $g_*$. This proves the invertible endpoint comparisons;
naturality for arbitrary endpoint squares is the factorization just
given.

\emph{Every middle balance.} For $u:B\to B'$ and $r:E\to B'$,
write $r_u:E_u\to B$ and $v:E_u\to E$ for its pullback. Factor the
square from $n$ to $u_*n$ through the display lifts, then bend over
$v$. The resulting fixed-target comparison is
\[
       v_* n_{r_u}\ \iso\ H^*(u_*n)_r .
 \tag{8}\label{eq:bc}
\]
Here $H:A_{r_u}\cong A_r'$ over $A$ is induced by $Tv$; it is
invertible because $T$ preserves this pullback and
$a_{u_*n}=Tu\,a_n$. To prove the \emph{cell} invertible, paste
$C_{r_u}(n)$ with the pure pullback square $\pure(v)\to\pure(u)$.
This is a cartesian square with bottom $n\pure(u)\cong u_*n$,
right side $r$, and top $n_{r_u}\pure(v)\cong v_*n_{r_u}$.
Comparing it to $C_r(u_*n)$ by cartesian uniqueness gives exactly
\eqref{eq:bc}: both are factorizations of the same square after
bending over $v$. Conversely, in \eqref{eq:bc} take $n$ a
horizontal identity. Its cartesian lift is the horizontal identity
on $r_u$; unbending the canonical isomorphism recovers the actual
pure pullback square as cartesian. This proves the asserted equivalence
in Definition~\ref{def:corner}. On decorations the balance comparison is
\[
     a_n^*\Tt(u^*P)\cong a_n^*(Tu)^*\Tt P
       \cong (Tu\,a_n)^*\Tt P .
\]
It lies over the same $H$. Append the target tail. The mate for a
composite of pure maps equals the composite of the mates with
whiskerings: restore both right cartesian arrows and unbend the
intermediate pure actions; the resulting pure pullback squares paste
to the pullback for the composite and the original input square is
the same. Uniqueness proves the equality; identity and naturality
with \eqref{eq:square} use the same argument. For the decorated part
this is uniqueness of restriction in $\D_1$. Thus \eqref{eq:cut}
descends through every defining relation of the pseudocoend.

\emph{Four comparisons.} Write $-P$ for the backward context with
identity tail. It suffices first to compute at consecutive content
arrows $n,m$ and consecutive backward displays $-P,-S$; how all
tails and balances enter is checked below. The following are the
generating components, showing the arrows up to inversion to match
the common direct-to-staged convention:
\begin{equation}\tag{8a}\label{eq:four}
\begin{array}{@{}lll@{}}
\text{source}& \text{decorated part}&\text{content comparison}\\
j_N M &U_{\id}(P):L_{\eta}P\cong P&
 C_r(1)\ \text{vs.\ the horizontal identity on }r\\
N j_M&I_{a_n}:L_{a_n}\un\cong\un&
 C_{\id}(n)\ \text{vs.\ the vertical identity on }n\\
m_N M&K_{a_n,a_m}(P)^{-1}&
 C_r(nm)\ \text{vs.\ horizontal paste }C_\ell(n),C_r(m)\\
N m_M&F_{a_n}(P,S)&
 C_{rs}(n)\ \text{vs.\ vertical paste }C_s(n_r),C_r(n).
\end{array}
\end{equation}
Here $\ell$ belongs to $C_r(m)$ in the third row.
For each row the two indicated cartesian squares have the same
right side and bottom, up to the indicated one-sided unitors and
associators. The unique comparison of their tops is invertible;
its left boundary is \emph{exactly} the carrier map in the middle
column. For example for horizontal composition its two projections
are the same original $A$-coordinate and the $\mu$-flattening
of the two Kleisli decorations, by \eqref{eq:arity}; the pullback for
the direct lift forces that map. For vertical composition the two
projections are the original $A$-coordinate and the last
$TF$-coordinate, with the intervening $TE$-coordinate obtained by
$Ts$. For the unit rows the two projections are forced by
cartesianness of $\eta$ and the identity pullback, respectively.
Thus each row gives a single invertible \eqref{eq:G} arrow.
As modifications these have types
\[
\begin{aligned}
 \tau(j_N M)&\cong Mj_N,&\qquad \tau(Nj_M)&\cong j_MN,\\
 \tau(m_N M)&\cong(Mm_N)(\tau N)(N\tau),&
 \tau(Nm_M)&\cong(m_M N)(M\tau)(\tau M).
\end{aligned}
\]

\emph{Extension from the displayed components.}
In the source $N^iM^j$ an internal tail $g$ of the $M$ block
can be removed by the \emph{invertible balance} between
$[(P,g),(S,f)]$ and $[(P,\id),(g^*S,fv)]$, where $v$ is the
pullback map into $\dom S$. The last tail is a pure output action.
The same argument removes the tail next to any further $N$ by
precomposition. These are isomorphisms of representatives, not
equations in a coend of sets. Naturality on product arrows follows
from naturality in Lemma~\ref{lem:local} and from factoring the
pasted square between lifts in \eqref{eq:square}.
For clarity, there are only three internal seams on which compatibility
with balance must additionally be checked.
\begin{enumerate}
\item At $N|N$, inserting a pure $u$ between two horizontal arrows
compares the two ways of cutting the same three-column horizontal
paste, with middle bottom $\pure(u)$. The central cartesian square
is the ordinary pure pullback by Definition~\ref{def:corner}(3);
the factor and its mate are \eqref{eq:bc}. The two grade maps
agree by the horizontal associator then cartesian uniqueness.
On decorations this is associativity of $K$ with middle arrow
$\eta u$, followed by $U_u$.
\item At $N|M$, insertion of $u$ compares a terminal pure horizontal
column and restriction of the first display along $u$. Its content
square is exactly the defining square of \eqref{eq:bc}.
For a vertical stack, the pure pullback of a composite display is
the vertical paste of its two pure pullbacks. For a horizontal
string the pure column pastes with the last content column.
The two factorizations into the terminal cartesian square are
therefore equal, including the original right pure mate by unbending.
On decorations the first assertion is horizontal compatibility of
$U_u$ with $F$, the second is the $K$ composition equation for a
final pure arrow (and cartesian preservation by $\Tt$).
\item At $M|M$ it is the slide $g^*S$ described above. Inserting
its pure square between two vertical display layers and pasting, or
first making the slide and pasting, gives the same large content
square by horizontal/vertical interchange with that pure pullback
square. Restoring the final cartesian square forces equality of the
factors. On decorations it is naturality of $F$ in this cartesian
restriction, with the ordinary identification $L_{\eta g}\cong g^*$.
\end{enumerate}
The endpoint versions are the first and last such slides, or postcomposition
of tails. Pasting two slides uses mate pasting as in \eqref{eq:bc};
slides in disjoint slots commute by functoriality. This proves that
the four rows give modifications on the actual relative tensors.
In particular a composition of unrestricted forward tails is not
treated as a display.

\emph{The full pastings.} The following cancellation observation
makes the grade check uniform. For a rectangle of consecutive content
columns and display rows, fill from bottom right with the chosen
cartesian squares. Vertical composition of cartesian arrows and the
horizontal stability hypothesis make every subrectangle cartesian.
Comparison to a single lift has the defining equation
$C\,t=\text{pasted rectangle}$, with the chosen bracketing and
bottom structural comparison. A comparison on a subrectangle,
whiskered and then compared on the whole rectangle, still satisfies
this equation by interchange and naturality of the horizontal
associator. Hence any two such composites, with the same
bottom one-sided map, are equal by cartesian uniqueness.
The converse orientations cause no problem, since the comparisons
are invertible. An empty row is a vertical identity and an empty
column a horizontal identity on the display. If source and target
bracketings differ, insert the one-sided associator/unitor on the
bottom and right; the one-sided pentagon and triangles identify
the restored boundaries. This argument compares actual content
squares before bending; combined with \eqref{eq:normal}, it works
even for nonfaithful arity. It does not establish the decorated
column of the following table; that column uses Lemma~\ref{lem:local}.

Use the numbering C1--C10 of \cite[Def.~4.1 and App.~A]{GL},
with $S=N$, $T_{\mathrm{GL}}=M$. Subscripts indicating restriction
and canonical bracket changes are suppressed in the elementary
identities below; all are identities on the common source representative.
\begin{center}\small
\begin{tabular}{@{}clp{.53\linewidth}@{}}
\toprule
 &Restored content diagram&Equality of the two decorated composites\\
\midrule
C1&three columns, one row&$K$ associativity, i.e.\ the two $\bar\mu$ multiplications on $\Tt^3P$.\\
C2&left unit column&$K$ left unit, the first monad unit identity.\\
C3&two columns, identity row&Horizontal-unit axiom of $\bar\mu$: flattening before or after $I$ on the horizontal unit.\\
C4&both units&Horizontal-unit axiom of $\bar\eta$: its restriction of the decorated unit is the chosen unit.\\
C5&two columns, two rows&The two sides of \eqref{eq:middle}, or inverses according to the chosen orientations.\\
C6&identity column, two rows&Horizontal-composition axiom of $\bar\eta$: $F$ at a pure unit agrees with ordinary composition of restrictions.\\
C7&one column, three rows&$F$ associativity on three consecutive displays, the compositor equation of $\Tt$.\\
C8&last identity row&$F$ right unit, compatible with the corresponding $\D$-unitor.\\
C9&right unit column&$K$ right unit, the other monad unit identity.\\
C10&first identity row&$F$ left unit, compatible with the other $\D$-unitor.\\
\bottomrule
\end{tabular}
\end{center}
Each table entry compares the components of its pasting on a
generating object with identity tails. By the preceding balance and
endpoint checks every other representative is connected to this
calculation by the indicated isomorphisms and output action.
Flattening an iterated pseudocoend, or rebracketing it, represents the
same multiply pseudobalanced cocone. Its associators thus make
exactly the slides already checked. Gray interchangers between
independent slots act as naturality squares, not extra coefficient
relations. Since these are modifications, equality of their
components on all objects is equality of modifications. The two
separate columns give equality in \eqref{eq:G}; the equivalence
\eqref{eq:normal} reflects it. This establishes all ten equations.
The last assertion now follows from composition of pseudomonads
with a pseudodistributive law, or the relative realization of this
composite in \cite[Sec.~6]{U}.
\end{proof}

\section{What naturality means}\label{sec:maps}
Here is a category of compatible inputs; allowing lax maps of all
possible sorts is unnecessary. A map from unprimed to primed data is:
\begin{enumerate}
\item a pullback-preserving functor $V:\B\to\B'$ taking $\J$ to $\J'$,
and a natural transformation $\theta:VT\to T'V$, cartesian on
displays, with
$\theta\,V\eta=\eta'V$ and
$\theta\,V\mu=\mu'V\,(T'\theta)(\theta T)$;
\item a pseudo-double functor $D:\D\to\D'$ over $V$ on maps and
carriers, preserving cod-cartesian squares, and a vertical double
transformation $q:D\Tt\to\overline{T'}D$ over $\theta$ whose
display squares are cod-cartesian. It obeys
$q\,D\bar\eta=\bar\eta'D$ and
$q\,D\bar\mu=\bar\mu'D\,(\overline{T'}q)(q\Tt)$;
\item a strong pseudofunctor on the content bicategories $G:N\to N'$
over $V$ on objects, equipped with a compatible invertible pure icon
$G\pure\cong\pure' V$. Its induced functor on bent squares preserves
the stipulated cod-cartesian squares and
$a'_{Gn}=\theta_B\,V a_n$.
\end{enumerate}
Composition uses the ordinary composites of the functors, whiskered
composites of the icons, and the pasted transformations $\theta,q$.
The defining equations are stable under whiskering and pasting;
composites of cartesian squares are cartesian. This gives a class
closed under identities and associative composition with the usual
data of strong functors (or a category after fixing bracket conventions).
A noninvertible $\theta$ is allowed; the necessary invertibility is
the induced \emph{display restriction} comparison, not invertibility
of $\theta_B$.

The induced context map sends $(P,g)$ to $(DP,Vg)$; its strength for
context multiplication is preservation of cartesian restriction followed
by the horizontal compositor of $D$. On relative tensors over different
bases, send the middle object $B$ to $VB$ and its balance along $u$
to the balance along $Vu$. The map of crossings compares
$(D_M\odot G)\tau$ to $\tau'(G\odot D_M)$ over $V$. Its decorated
component is
\[
 D(a_n^*\Tt P)\ \cong\ (V a_n)^*(D\Tt P)
 \ \cong\ (V a_n)^*\theta_B^*(\overline{T'}DP)
 \ \cong\ L'_{a'_{Gn}}DP .
\]
Its other component is the comparison of $G C_r(n)$ and $C'_{Vr}(Gn)$,
followed by pure transport over $Vg$. They have the same boundary
because the arity of the former projects by $\theta_E$ and both are the
same target display pullback. Naturality on \eqref{eq:square} and
\eqref{eq:bc} follows by restoring the cartesian squares; the pure
icon preserves the opcartesian factors. For \eqref{eq:four}, the
decoration equations are, respectively, $q$ preserving the monad
unit, the horizontal functor unit, $q$ preserving the monad
multiplication, and $q$ preserving the horizontal compositor.
On content the restored comparisons are the same images of the
one-sided horizontal unit, vertical identity, horizontal and vertical
pastes; uniqueness identifies their factors. For two composable input
maps the two decoration comparisons restrict the pasted $q$; the two
content comparisons have the same restored cartesian projection.
Consequently the comparison for a composite is the pasted comparison,
and that for an identity is the identity with the stipulated
cleavages. This proves parameter naturality of the law and the induced
strong map on its composite. No map of mixed laws belongs to the input.

\section{Building the lifted monad by finite polynomial evaluation}\label{sec:rig}
Let $\B=\Set$ and $TX=\coprod_{s\in S}X^{F_s}$ be a cartesian
polynomial monad with finite addressed $F_s$. Let $\V$ be a symmetric
rig category: both $(\oplus,0)$ and $(\otimes,I)$ are symmetric monoidal,
with the coherent invertible finite distribution and absorption maps
of a distributive rig. Choices of indexed brackets and bijections are
fixed with their canonical structural maps. Neither monoidal operation
need be categorical (co)product. Use decorated displays of finite-fibre
$r:E\to B$ in $\J$; if $\J$ contains maps with infinite fibres,
the decorated maps form its finite-fibre part, which is closed under the
required operations. A decoration at $b\in B$ is $(V_e)_{e\in E_b}$
and \emph{one} arrow $\epsilon_b:\bigoplus_{E_b} V_e\to I$.
This is not a family of arrows $V_e\to I$.
Over a square $(h,u)$ a map has one coefficient map for each
$b\in B$ and $e'\in E'_{u b}$:
\[
 \gamma_{b,e'}:\bigoplus_{\substack{e\in E_b\\ he=e'}}V_e
                      \longrightarrow V'_{e'}.
 \tag{9}\label{eq:nat}
\]
After repartition by $h$, their sum followed by $\epsilon'_{ub}$
must equal $\epsilon_b$. In particular fibres are taken over the
\emph{given} $b$ even when $u$ merges base points. Composition
partitions the fibre of a composite before applying the actual
coefficient arrows. The restriction along $u$ uses the ordinary
pullback and identical fibres; the unique factor of a square uses
the same coefficient maps over the pullback's fibre bijection.
This proves the full cod-cartesian property.

For a finite set $X$, let $S_XV=\bigoplus_X V_x$. A useful derivation
of \eqref{eq:nat} starts with finite bispan evaluation
\[
 (L\xleftarrow{s}U\xrightarrow{p}D\xrightarrow{t}K)
 \quad\mapsto\quad
 \left(V\mapsto
   \left(\bigoplus_{td=k}\ \bigotimes_{p i=d} V_{s i}\right)_{k\in K}\right).
 \tag{10}\label{eq:bisp}
\]
It acts on tuples of arbitrary $\V$-arrows. Under substitution,
distribution indexes each summand by one choice per input position;
the remaining positions are the disjoint union of their addressed
leaves. Comparing the same choices in a three-level substitution gives
the compositor equation by rig coherence. An isomorphism of diagrams
just reindexes summands and positions. This is the ordinary finite
bispan model relevant to \cite[Thm.~A]{CHLL}, evaluated fibrewise
when $B$ is infinite. To obtain directed carrier maps, \emph{restrict}
to covariant additive transfers, $h_!V=(\bigoplus_{hx=y}V_x)_y$,
and form their Grothendieck category $\mathcal R$ with arrows
$(h,\gamma:h_!V\to V')$. Total sum defines
$S:\mathcal R\to\V$, using fibre partition. The augmented category is
$\A=\mathcal R\times_\V \V_{/I}$ along the domain functor.
This proves \eqref{eq:nat} and its composition; it does not infer
directed arrows from the $(2,1)$-source of \eqref{eq:bisp}.
For example in the discrete rig of natural numbers, a singleton
augmented fibre has one object, while a two-element carrier has two,
$(1,0),(0,1)$; augmented fibres do not preserve the disjoint-union
products of the raw model.

\paragraph{The horizontal operations, on their arrows.}
For $S':F\to E$ over $s$ with weights $W_z$ and augmentations
$\delta_e$, the composite $P\st S'$ has carrier $rs$, row
$V_e\otimes W_z$ at $z$ with $sz=e$, and augmentation obtained by
partition by $e$, distribution, the $\delta_e$, and finally $\epsilon_b$.
For outer maps over $h$ with coefficients $\alpha$ and inner maps
over $k$ with coefficients $\beta$, the coefficient at $b,z'$
($s'z'=e'$) is the following map; sums are in the source fibre at $b$:
\begin{align}
 \bigoplus_{he=e'}\ \bigoplus_{\substack{sz=e\\kz=z'}}
                 V_e\otimes W_z
 &\cong\bigoplus_{he=e'} V_e\otimes
                      \big(\bigoplus_{\substack{sz=e\\kz=z'}}W_z\big)
 \xrightarrow{\oplus_e(1\otimes\beta_{e,z'})}
        \bigoplus_{he=e'} V_e\otimes W'_{z'} \notag\\[-2pt]
 &\cong\big(\bigoplus_{he=e'}V_e\big)\otimes W'_{z'}
 \xrightarrow{\alpha_{b,e'}\otimes1}V'_{e'}\otimes W'_{z'} .
 \tag{11}\label{eq:dep}
\end{align}
Continuation maps at different ancestors need not agree; all are
applied before the outer merge. The subsequent factor $W'_{z'}$
\emph{is} constant on that merge. Under a further square, after
partition of the composite fibre, the first outer map commutes
with the subsequent inner map because the latter is fixed on each
already merged intermediate $e'$ and tensor is bifunctorial.
Distributing before and after this exchange gives the composite
formula \eqref{eq:dep}. This proves interchange on arbitrary maps;
the two augmentation equations in inside-first order prove that the
result is augmented. The unit is the singleton $(I,\id)$;
rebracketing triple composites on an addressed chain gives the tensor
associator, and augmentation contracts the same last-to-first chain.
The same constant-on-merged-fibre exchange proves its naturality;
the monoidal pentagon and unit coherence give the pseudo-double
category axioms. Cartesian restrictions paste horizontally since the
two constituent pullback fibre maps are bijective with invertible row
maps, as is their composite, with the same augmentations. More
explicitly, a square of decorated displays over a pullback is
cod-cartesian when its coefficient maps under the fibre bijections
are isomorphisms: it is the canonical cartesian restriction preceded
by a vertical isomorphism. Conversely cartesianness forces that
isomorphism.

\paragraph{Occurrence lift, its action on maps, and its structural squares.}
Over $t=(s,(b_i)_{i\in F_s})\in TB$ the display $\Tt P$ is
$Tr:TE\to TB$ with fibre $\prod_iE_{b_i}$; at $(e_i)$ its weight
is $\bigotimes_i V_{e_i}$. Its total augmentation is
\[
 \bigoplus_{(e_i)}\bigotimes_i V_{e_i}
 \cong \bigotimes_i\bigoplus_{e_i} V_{e_i}
 \xrightarrow{\otimes_i\epsilon_{b_i}}\bigotimes_i I\cong I.
\]
A square with coefficients $\gamma$ gives, at a target tuple $(e'_i)$,
\[
 \bigoplus_{\forall i:\, h e_i=e'_i}\bigotimes_i V_{e_i}
 \cong \bigotimes_i\bigoplus_{h e_i=e'_i}V_{e_i}
 \xrightarrow{\otimes_i\gamma_{b_i,e'_i}}\bigotimes_i V'_{e'_i}.
 \tag{12}\label{eq:theta}
\]
These sums too are in their specified $b_i$-fibres. This is an inverse
distributor (comparison of two polynomial evaluations), followed by
the actual directed maps. For two squares the two routes partition
the fibre first by the intermediate tuple, then by its coordinates;
naturality of distribution gives precisely the same tensor of
composites. Totalling \eqref{eq:theta} and using each incoming
augmentation equation gives the outgoing augmentation equation.
Pullback restriction is preserved since row bijections and isomorphisms
remain such.

The comparison $\Tt(P\st S')\cong\Tt P\st\Tt S'$ splits each tuple
of dependent pairs into its outer tuple and its dependent inner tuple;
on coefficients it shuffles
$\bigotimes_i(V_i\otimes W_i)\cong
(\bigotimes_i V_i)\otimes(\bigotimes_i W_i)$.
At a target tuple $(z'_i)$, insert \eqref{eq:dep} in
\eqref{eq:theta}: every summand records a \emph{single} compatible
pair $(e_i,z_i)$ at each occurrence. Splitting the tuple first gives
the same sum over outer tuples $(e_i)$ followed by
$z_i\in F_{e_i}$ with $kz_i=z'_i$. Distribute the sums by $i$;
on both routes apply $\beta_{e_i,z'_i}$ before $\alpha_{b_i,e'_i}$.
All other exchanges commute by tensor bifunctoriality on disjoint
factors or by distributor naturality at a common target row. Hence
the shuffle is natural on the actual aggregations, not only on their
supports. Both augmentations first use the inner augmentations and
then the outer ones, up to exchanges on independent tensor factors.
The horizontal unit comparison collapses the forced $I$ factors.
For three displays both routes split the same tuple of ancestor
chains $(e_i,z_i,w_i)$, and induce the same block permutation and
rebracketing; this proves the horizontal compositor equation.

A cartesian polynomial unit has one position. A cartesian polynomial
multiplication has a specified leaf bijection
$F_{\mathrm{flat}(s,(s_i))}\cong\coprod_i F_{s_i}$.
Both statements follow by evaluating the position map of a natural
transformation between sums of representables on its generic labelling;
cartesianness makes that map a bijection. Define
$\bar\eta_P,\bar\mu_P$ on the carriers by $\eta,\mu$ and on rows by
the singleton identification and addressed tensor flattening,
respectively. Underlying pullbacks and invertible row comparisons
make them cod-cartesian. Their naturality on \eqref{eq:theta} is
partition by the same leaves, with the same tensor of actual
aggregations. For the nontrivial horizontal transformation axiom of
$\bar\mu$, both routes send the two-level family of pairs
$(V_{ji},W_{ji})$ to the block of all $V_{ji}$ and the block of all
$W_{ji}$ in the same leaf order. The monad equations identify the
bijections of leaves on three nested polynomials, not just their
cardinalities: equality of the two polynomial transformations
evaluated on the generic labelling gives equality of their position
maps. Tensor coherence now proves the monad and unit equations as
decorated arrows. This constructs Definition~\ref{def:lift} without
invoking any content object.

For precision about coherence in these calculations, tag primitive
occurrences before identifying variables with equal values. The
distributive expansions used above index monomials by compatible
histories, not by equality of their interpreted factors. In comparing
structural reindexings, use the indicated ordering/bijection of those
histories and of their leaves. In the usual binary rig presentation
this is regular distributive coherence \cite{Lap}: distinguish
independent inputs and tag nullary nonzero alternatives by distinct
unit variables before specializing them to $I$. Repeated appearances
of a factor across alternative summands do not repeat it \emph{in}
a monomial; for example a common target row in \eqref{eq:dep} is used
only after the corresponding alternatives have merged. Naturality
moves nonstructural arrows through the structural maps where the
incidences match. Empty sums and products use absorption and unit
coherence, not uniqueness of arrows from $0$ or into $I$.

\section{Recovering the independent Set law, including its cells}\label{sec:set}
Here is a content-only criterion for Definition~\ref{def:corner}; its
proof also identifies the exact grades in the specialization. It is
stated over $\B$ because the next example uses it there. Take a valid
one-sided comma content presentation $p:\Lcat\to\B$, $d:\Lcat\to\B$:
$N(A,B)=(A\downarrow d_B)$, with $Q\in\Lcat_B$, $f:A\to dQ$;
unit $U_B$ with body $\eta^L_B:B\to dU_B$, and dependent accumulation
$Q\diamond R\in\Lcat_B$ for $R\in\Lcat_{dQ}$ with body
$\mu^L_{Q,R}:dR\to d(Q\diamond R)$. These need not be invertible.
A total grade square $m:Q\to Q'$ over $u$ has $dm\,f=f'k$.
Composition of $(Q,f)$ followed by $(R,b)$ has grade
$R\diamond b_!Q$ and body $\mu^L_{R,b_!Q}\,d\kappa_{b,Q}\,f$,
where $\kappa$ is the chosen opcartesian transport for the global
opfibration $p$. Assume precisely the following additional data here:
\begin{enumerate}
\item $p$ has genuine cartesian lifts $c_{r,Q}:r^*Q\to Q$ over $\J$,
tested against arbitrary total arrows. There is $\nu:d\Rightarrow Tp$
natural on all grade arrows with
$\nu_{U_B}\eta^L_B=\eta_B$ and
$\nu_{Q\diamond R}\mu^L_{Q,R}=\mu_B(T\nu_Q)\nu_R$.
\item $(dc_{r,Q},\nu_{r^*Q}):d(r^*Q)\xrightarrow{\sim}
dQ\times_{TB}TE$ is the canonical display isomorphism. Call it
$\Delta_{r,Q}$. Thus $j=dc_{r,Q}$ lies in $\J$.
\item $U$ sends displays to cartesian arrows. Accumulation sends the
cartesian pair $(c_{r,Q},c_{j,R})$ to a cartesian arrow above $r$.
The canonical mate $v_!r_u^*Q\to r^*u_!Q$ of every pure pullback square
in Definition~\ref{def:corner}(3) is invertible. All four edges of the
square exist as pure maps, not necessarily as display changes of grade.
\end{enumerate}

\begin{lemma}[Actual coextension]\label{lem:comma}
These conditions give relative corners with $a_{(Q,f)}=\nu_Q f$.
\end{lemma}
\begin{proof}
Naturality and the two equations for $\nu$ prove \eqref{eq:arity},
using naturality at $\kappa_b$ in the composite. Against $r:E\to B$
the chosen lift has grade $r^*Q$ and source
$A_r=A\times_{dQ}d(r^*Q)$ with its second projection as body.
A testing square with grade $l:Z\to Q$ over $rv$ factors uniquely
in grades through $c_{r,Q}$ over unrestricted $v$. Its unique left
edge is then forced by the pullback using the factored grade and the
given source body. This proves cartesianness in $\N_1$, not just in
$\Lcat$. Display identifies the arity square with the required
pullback. For the pure arrow $u:B\to B'$ the lift along $r:E\to B'$
has grade $r^*U_{B'}\cong U_E$ and source
$B\times_{TB'} TE\cong B\times_{B'}E$ by cartesianness of $\eta$.
The top body is $\eta^L_E v$ since its two display coordinates are
$(\eta^L_{B'}u r_u,\eta_E v)$. This proves pure exactness on the
\emph{canonical} pure square.

For the horizontal paste let $(Q,f):A\rightsquigarrow B$ be followed
by $(R,b):B\rightsquigarrow C$ and lift against $r:E\to C$.
Write $R'=r^*R$, $j=dc_{r,R}$,
$B'=B\times_{dR}dR'$ with projections $v:B'\to B$, $b':B'\to dR'$.
The pasted grade, with its original projection, is
\[
 R'\diamond b'_!v^*Q\ \cong\ R'\diamond j^*b_!Q
                \ \cong\ r^*(R\diamond b_!Q).
\]
The first comparison is the mate for $(v,j;b,b')$, the second the
cartesian accumulation comparison. Their factorization equations
identify the pasted \emph{grade arrow} with the cartesian lift of the
actual composite grade. For its source let $a=\nu_Q f$, $c=\nu_R b$.
The two-step pullback is
\[
 A\times_{TB}TB'\ \cong\ A\times_{TTC}TTE
              \ \cong\ A\times_{TC}TE .
\]
Here $B'\cong B\times_{TC}TE$; the first comparison uses $T$ preserving
pullbacks, over $(Tc)a$, and the second uses the cartesian square
of $\mu$ at $r$. The bottom composite has arity $\mu_C(Tc)a$,
and the top composite has the corresponding flattened arity by
\eqref{eq:arity}. Its other display coordinate is its original
composite body, since the pasted grade/body square is an actual
content square. Thus this comparison identifies the top arrow with
the direct lift, not just its source. This proves horizontal stability
including closure of all intervening left legs. No preservation of
opcartesian arrows by $d$ was used.
\end{proof}

\begin{corollary}[Recovery with the original marks]\label{cor:recover}
Under the Set input of \cite{O1,C} (symmetric rig, finitary
cartesian polynomial monad, valid relative comma content and the
three clauses above), Theorem~\ref{thm:glue} gives the same crossing,
balances, constraints, and composite on the original hom categories as
the independent law. In particular its new content grade is the
\emph{actual} relative pullback of the old grade before forward
transport, without restriction to the image of a body.
\end{corollary}
\begin{proof}
Use \eqref{eq:normal} with comma $N$; then $\grade_N(A,C)$ is
literally (up to the chosen pseudofunctor actions)
\[
 H(A,C)=(\ea_A\downarrow d_C).
\]
Its objects are $(P,Q,f:\ea_AP\to dQ)$; its arrows are native
$(\alpha,\beta)$ with $d\beta\,f=f'\ea\alpha$. Over endpoint
maps use the same equality after the opcartesian target and
cartesian source transports. For a display $\pi:E\to B$ of $P$
and $Q$ over $B$, put $\widehat P=L_{\nu_Q}P$,
$\widehat Q=\pi^*Q$. At $y$ with $\nu_Q(y)=(s,(b_i))$,
$\widehat P_y=\bx_{i\in F_s}P_{b_i}$ with precisely the total
augmentation and native arrows of \eqref{eq:theta}. Its carrier is
$dQ\times_{TB}TE$, and its generic body is $\chi=\Delta_{\pi,Q}^{-1}$
after this carrier identification. On $(Q,f:A\to dQ),(P,g:E\to C)$
the resulting comma object is
\[
 \left(f^*\widehat P,\ g_!\pi^*Q,\quad
      \ea(f^*\widehat P)\to\ea\widehat P
       \xrightarrow{\chi}d(\pi^*Q)
       \xrightarrow{d\kappa_{g,\pi^*Q}}d(g_!\pi^*Q)\right).
 \tag{13}\label{eq:setcross}
\]
For $m:Q\to Q'$ over $u$ and a decorated square with carrier $h$ the
native grade comparison is the opcartesian factor of the unique
$\widetilde m$ over $h$ with
$c_{\pi',Q'}\widetilde m=m c_{\pi,Q}$:
\[
              h_!\pi^*Q\longrightarrow\pi'^*Q'.
\]
It is not assumed an isomorphism. The context map is the distributor in \eqref{eq:theta}
followed by the given aggregates there, not a function of $h$
alone. The body equation holds because its target display coordinates
are $(dm,Th)$ applied to the old coordinates, by the factorization
equation and naturality of $\nu$. For a middle pure map $u$, the
grade comparison is exactly $v_!\pi_u^*Q\cong\pi'^*u_!Q$ by the
same defining equations, and the corresponding two body coordinates
are $(d\kappa_{u,Q}\,y,Tv(t))$. This is the chosen whole-base balance.

On \eqref{eq:four}, the first row uses the one-position identification
\emph{after restriction along} $\eta^L_B$ and $U_E\cong\pi^*U_B$.
The second uses forced singleton choices and $(\id_B)^*Q\cong Q$.
The third uses addressed leaf regrouping \emph{after restriction along}
$\mu^L_{Q,R}$ and the inverse, or direct comparison depending on
orientation, of
$\pi^*Q\diamond(dc_{\pi,Q})^*R\cong\pi^*(Q\diamond R)$.
In particular the first display coordinate is $\mu^L_{Q,R}(z)$,
not $z$. The fourth uses the actual shuffle natural on
\eqref{eq:dep} and $(\pi\rho)^*Q\cong\rho^*\pi^*Q$. The grade
isomorphisms preserve their original projections by construction.
These are the defining components of the independent interaction
\cite[Secs.~4--5]{O1}; all product arrows, balance generators, and
four modifications agree after the indicated canonical reindexings.
The ten pastings in Theorem~\ref{thm:glue} therefore specialize on
the same cells. For successive $(P,Q,f),(D,S,g)$ the output context
is $P\st f^*\widehat D$ and content is $S\diamond g_!\pi_D^*Q$.
This identifies the chosen multiplication as well as the category
on which it acts.
\end{proof}

\section{A genuinely indexed test: restriction erases a colour and merges rows}
\label{sec:diagram}
Let $\B$ be the presheaf category on $0\to1$. Write an object as
$X_1\to X_0$, with \emph{every} natural transformation a pure map.
Take pointwise finite-fibre maps for $\J$. Fix a nontrivial group
$G$ (two elements suffice). Set
\[
 (TX)_1=\coprod_{n\ge0}(G\times X_1)^n,\qquad
 (TX)_0=\coprod_{n\ge0} X_0^n.
\]
Its restriction erases each $G$-label and applies $X_1\to X_0$ to
the entry. The unit is a singleton with label $1_G$ at stage 1.
Flatten nested words there by concatenating and multiplying the
outer colour with each inner colour in that order. Unit and
associativity follow from concatenation and the group laws, and the
restriction commutes with both. This is an internal finite-position
polynomial: the shape presheaf is
$\coprod_nG^n\to\mathbb N$, the map returning length; positions are
pairs of a shape and an index $1\le i\le n$, restricting to the
same index. Over a generalized shape $s:y c\to\text{shapes}$ the
position object is a coproduct of $n$ copies of the representable
$yc$; a compatible labelling is precisely $n$ elements of $X_c$.
Thus the polynomial really has the displayed stagewise formula.
The unit and multiplication have bijective addressed position maps
at each stage, hence are cartesian; $T$ preserves pullbacks
stagewise and $T\J\subseteq\J$. The polynomial is not the constant
list polynomial pulled from $\Set$: at length one the shapes restrict
by $G\to1$.

Use the same rig at both stages. An internal decorated display is a
diagram of decorated finite-fibre Set displays: in addition to its
stage objects it includes a square of the form \eqref{eq:nat} over
$(E_1\to E_0,B_1\to B_0)$, not necessarily invertible or cartesian.
A map is a natural square of such diagrams: the square must commute
on \emph{coefficient arrows}. Horizontal units and composition are
pointwise, with the restriction square computed by \eqref{eq:dep}.
To restrict along a natural transformation, take both pointwise
cartesian restrictions and use the universal property of the lower
one to induce the comparison on the restriction arrow. A factoring
square of diagrams factors uniquely at each stage; naturality of
the factors follows by restoring the lower cartesian square.
Thus the chosen squares are genuinely cartesian against all diagram
squares. Their underlying squares are pointwise pullbacks, and they
are stable under horizontal pasting.

Lift $T$ using the coloured occurrence product at stage 1 and its
uncoloured version at stage 0. On the restriction arrow use
\eqref{eq:theta} in the \emph{same} positions, forgetting the
colours, with the given restriction aggregations. Identity and
composition of maps of diagrams follow from the composite-fibre
equation for \eqref{eq:theta}; in particular the restriction map
is not made into a formal isomorphism. The horizontal compositor
is a map of diagrams by the naturality of the addressed shuffle on
\eqref{eq:dep}. The squares for $\bar\eta,\bar\mu$ are maps of
diagrams because forgetting colours preserves the specified leaf
bijections and the map for an aggregation under flattening is the
same tensor of its leaf maps. Their cartesianness follows by the
pointwise factorization argument above, as does preservation of
cartesian squares. All remaining equalities of maps of diagrams
are detected jointly by stage evaluation \emph{after} this naturality
check. This verifies the independent decoration input without assuming
that pointwise finite fibres are locally constant internal finite sets;
they need not be. It is the \emph{positions of the polynomial} which
restrict bijectively, not the rows of every display.

For independent content take receipts $\Lcat_B=\B/(TB)$. A grade is a
tag presheaf $q:Q\to TB$, pushforward postcomposes by $Tu$, and
$r^*$ is pullback of tags along $Tr$. Define
$dQ=B+Q$, $\nu_Q=[\eta_B,q]$, $U_B=(0\to TB)$,
$\eta^L$ the coproduct-unit identification. For $t:R\to T(B+Q)$,
the accumulated tags are $Q+R$ with label
\[
 [q,\,\mu_B(T\nu_Q)t] : Q+R\to TB;
\]
the body is coproduct reassociation. A map acts by its actual tag
maps on the two summands. These are natural presheaf maps. The
one-sided unit leaves all preexisting tags unchanged; the labels
of a third tag stage on either association are equal by $\mu$
associativity and naturality; the original and second tag stages
are retained with the same labels. The tag-map equations and the
one-sided pentagon and unit laws follow from coproduct coherence
on these same summands, including pushforward, which only relabels.
This is not the false slice-domain one-sided algebra.

The canonical mate in a display pullback is invertible since $T$
preserves the pullback and pushforward keeps the entire tag presheaf.
Extensivity and cartesianness of $\eta$ give
\[
 d(r^*Q)=E+(Q\times_{TB}TE)
                  \ \cong\ (B+Q)\times_{TB}TE .
\]
The unit is cartesian. For accumulation the $Q$ summand restricts
as displayed; for the $R$ summand the two comparisons are
$T(d(r^*Q))\cong T(dQ)\times_{TTB}TTE$ and the cartesian
naturality square of $\mu$ at $r$. Pulling $R$ along the resulting
comparison identifies its restricted nested labels with the pullback
of its flattened label, \emph{over the same original $R$}. This proves
preservation of the original cartesian projection and hence cartesian
accumulation. All computations are natural in presheaves since they
use canonical pullbacks, coproducts and the internal monad maps.
Lemma~\ref{lem:comma} therefore gives the independent content corners.

\paragraph{One cut that is not a pair of unrelated Set cuts.}
Let $B$ be terminal; let $E_1=\{1,2\}\to E_0=\{*\}$ be the fold,
with display $r:E\to B$. Take $\V$ to be finite-dimensional rational
vector spaces with direct sum and tensor. At stage 1 use one copy
of $\mathbb Q$ on each row, total augmentation $(a,b)\mapsto a+b$.
At stage 0 use $W=\mathbb Q^2$ on the one row, with the same sum
augmentation. Two possible restriction maps $\mathbb Q\oplus\mathbb Q
\to W$ are $\id$ and $p(a,b)=(a+b,0)$; both are lawful because they
preserve the global augmentation. The resulting objects cannot be
identified by an invertible change of row coordinates: their
restriction ranks are different.

Take one receipt at each stage, restricting to the other receipt,
labelled at stage 1 by $[(g,*),(1_G,*)]$ for $g\ne1_G$, and at
stage 0 by $[*,*]$. Use the generic input body $\id_{dQ}$ and
context tail $\id_E$. At the receipt, the new grade $r^*Q$ has
four tags upstairs, the choices $(i,j)\in\{1,2\}^2$, and one tag
downstairs. Its presheaf restriction is the actual map $4\to1$,
not four independently surviving tags. The new context has four
one-dimensional upstairs rows and one downstairs row $W\otimes W$.
On this fibre \eqref{eq:theta} makes the structural restriction the
identification $\mathbb Q^4\cong W\otimes W$ followed respectively
by $\id$ or $p\otimes p$. Their ranks are $4$ and $1$; specifically
the $(2,2)$ row goes to $e_2\otimes e_2$ in the former and to
$e_1\otimes e_1$ in the latter. Both have precisely the same tag
fold and coloured-to-uncoloured label restriction. Thus the difference
is in actual internal reindexing, invisible to the pair of stage
object lists. There is no presheaf isomorphism identifying these
crossed contexts even if arbitrary invertible row changes are allowed.
This example tests the hypotheses outside the Set theorem and supplies
a noninvertible internal transition at the cut itself. If the same fold
and coefficient maps are instead used as native Set maps at a fixed
middle point, with one tag labelled by the ordinary length-two word,
the grade arrow \eqref{eq:square} before the forward tail is the
actual four-to-one tag map $h_!r^*Q\to(r')^*Q$, $r'=\id_1$.
Thus the concrete comparison also exhibits a noninvertible
\emph{native grade cell}, distinct from the presheaf restriction
just calculated.

\paragraph{A noninvertible parameter change.}
Let $T'$ be ordinary lists at \emph{both} stages. Forget colours at
stage 1 and act identically at stage 0, giving a monad map
$\theta:T\to T'$ on the same presheaf base. Its naturality squares
are pullbacks: for a prescribed coloured shape above a word, the
labels of a proposed uncoloured lift determine exactly one lift with
that colouring. The unit and multiplication equations were verified
by colour erasure under concatenation. In Section~\ref{sec:maps}
take $V,D$ to be identities; $q:\Tt\to\overline{T'}$ uses the
identical row tensor over that bijection of fibres. It is cartesian
and preserves the horizontal shuffle, unit and flattening.
On receipts $G$ postcomposes each tag label with $\theta_B$, leaving
tags and bodies alone. The monad map equations identify accumulated
labels and the pullback squares of $\theta$ identify restricted
tags, over their original projections. This gives a genuine input
map, not an equivalence: two differently coloured length-one shapes
can have the same image. Its comparison of cuts preserves the same
four upstairs choices and their actual linear aggregation while
forgetting a colour. Thus naturality here does not mean that the
parameter functors themselves are invertible.

\paragraph{Bounded replay.}
The accompanying standard-library script
\texttt{check\_gluing\_witnesses.py} compares the two projections of
the colour-erasure pullback for the binary carrier fold on words of
length at most three (85 elements), and erasure with substitution on
1,807 nested words with both lengths at most two. The two routes of
\eqref{eq:middle} on 7,483 coloured nested words with dependent
pairs $(0,0),(1,0),(1,1)$ agree on every addressed choice.
With ancestor-dependent continuation columns $e_2,e_1,
(e_1+e_2)/2$, the native two-position Fubini routes agree on all
nine basis inputs and the resulting nine columns have rank nine.
The distinct rank-four and rank-one restrictions above are also
checked by rational row reduction. These are regression witnesses,
not a finite verification of Theorem~\ref{thm:glue}; none is needed
to read or compile the proof.

\section{Scope, separating obstructions, and recognition}\label{sec:boundary}
\paragraph{Two different index sets.}
Even a finite pointwise display is not automatically an allowable
\emph{set of multiplicative positions} with arbitrary restriction
arrows. In the discrete rig $(\mathbb N,+,\cdot)$ let a singleton
row of weight $1$ restrict to the first of two rows of weights
$(1,0)$. This is a valid augmented transfer: the fibres require
$1=1$ and $0=0$. A putative natural operation tensoring \emph{all
witness rows}, however, would require an arrow $1\to0$ after
restriction, which does not exist. In Section~\ref{sec:diagram}
multiplication ranges over polynomial positions, which restrict
bijectively for a fixed shape; arbitrary witness-row changes use
\eqref{eq:theta} separately in each such position. This witness rules
out replacing Definition~\ref{def:lift} by pointwise finiteness
alone, not more sophisticated indexed norm doctrines with genuinely
additional structure.

\paragraph{Why relative is not global.}
A failure outside $\J$ need not affect the law. For instance on
$\Set$, take constant discrete grades the two endofunctors
$\id,G_0$, where $G_0(\varnothing)=\varnothing$ and
$G_0(X)=1$ otherwise. Composition of endofunctors supplies the
one-sided content algebra, with extension by evaluation; both grades
close under this composition. Give the identity the singleton list
arity, and $G_0$ the empty list arity. These assignments respect
unit and substitution even on the empty set. For surjections
$E\to B$ the two sets have the same inhabitation status; display
for $G_0$ holds since the unique list over the empty list is empty.
Unit, cartesian grade accumulation and mates hold for the constant
discrete grading. For $\varnothing\to1$ display on $G_0$ fails.
More sharply, let $i_0,i_1:1\to2$ have disjoint images. Follow the
$G_0$ arrow $1\rightsquigarrow1$ by pure $i_0$. Its composite has
source $1$ and $G_0$ grade at $2$. Lifting the composite along
$i_1$ gives $1\times_{G_0(2)}G_0(1)=1$; lifting $i_0$ first gives
left source $\varnothing$ and then lifting the $G_0$ arrow gives
$\varnothing$. Horizontal cartesian stability fails outside the
display class. One may take finite-fibre surjections for $\J$, but
both $i_0$ and $i_1$ still act as \emph{pure maps} in the relative
tensor. This test, from \cite{C}, distinguishes restricting lift
existence from restricting the vertical category.

\paragraph{Marked necessity is narrower.}
In the comma setting, let $\mathcal D\subseteq\J$ be the actually
realized context projections. Suppose an otherwise independently
given law is marked by the complete carrier $dQ\times_{TB}TE$, the
chosen projection $c_{\pi,Q}$, an invertible body $\chi$ preserving
both coordinates, and by the original unit, accumulation and pure
transport projections. Then $\Delta_{\pi,Q}\chi=1$ and
$\Delta_{\pi,Q}=\chi^{-1}$. Preservation of the unit and
accumulation projections says that their grade comparisons factor
the original $U(\pi)$ and the accumulation of
$(c_{\pi,Q},c_{dc_{\pi,Q},R})$ through the respective cartesian
arrows. Cartesian uniqueness forces these canonical factors;
invertibility says exactly that the original projections are
cartesian. Analogously the pure-middle projection mark forces
the canonical cartesian--opcartesian mate, so it must be invertible
on each tested pure pullback. Without invertibility of $\chi$ only a
section of $\Delta_{\pi,Q}$ follows. Without realization of a
display there is no test of it. These are marked necessities
\cite[Sec.~6]{O1}, not a reconstruction of the lifted monad from
any unmarked abstract law.

\paragraph{Sources and actual increment.}
Buckley--Garner \cite[Secs.~3,6, esp.\ Prop.~6.3]{BG} give the
cartesian-square crossing of formally composable arrows for a globally
co-crossed double category. Definition~\ref{def:corner} restricts
where squares are lifted without changing what tests their universal
property. Their product of formally composable arrows is not our
relative pseudocoend; the pure pullback condition and \eqref{eq:bc}
are precisely where that distinction enters. Cnossen--Haugseng--Lenz--Linskens
\cite[Thm.~A]{CHLL} identify commutative semirings in a cartesian
closed presentable infinity-category with product-preserving functors
out of finite bispans. Here the finite operator calculus supplies
the raw coefficient evaluation, not the new grade, global augmentation,
or dependent native arrow. Capucci--Myers \cite[App.~D, Def.~D.9]{CM}
include an intertwiner and its structure cells as wreath data over a
pseudocategory. The gluing theorem derives its four comparisons
from inputs on \emph{separate} sides, and descends to the mixed
category-valued tensor. It is not an assertion that these three
published ambient categories are equivalent.

The earlier comparison \cite{C} verifies the concrete Set cut by
identifying it with the already established independent law. Here
the theorem is proved \emph{before} such a law is available, for
arbitrary compatible relative content, an indexed occurrence lift,
all pure maps, and a specified class of maps between possibly
different bases; the presheaf cut is a second application.
Proposition~\ref{prop:matched}, cartesian uniqueness and the pseudocoend each supply a local
universal property. Neither is an initiality theorem for all possible
assembler data or all ways to combine context and content. Such a
claim requires a category of realizations, comparison morphisms
and an initiality or representation proof, none of which follows
from \eqref{eq:cut}. No claim about arity-changing native content
cells, arbitrary internal presheaf position maps, infinite tensor
products, lax/noninvertible distributors, or higher-categorical
bispan enhancement is made.

\begin{thebibliography}{9}
\bibitem{BG} M. Buckley and R. Garner, \emph{Composing Double Categories},
\href{https://arxiv.org/html/2610.41551v2}{arXiv:2610.41551v2},
Secs.~3 and 6.
\bibitem{CHLL} B. Cnossen, R. Haugseng, T. Lenz and S. Linskens,
\emph{Homotopical commutative rings and bispans},
\href{https://arxiv.org/html/2403.06911v2}{arXiv:2403.06911v2},
Theorem~A.
\bibitem{CM} M. Capucci and D. J. Myers, \emph{Contextful Arrows via
Wreath Products}, \href{https://arxiv.org/pdf/2410.21889v2}{arXiv:2410.21889v2},
Appendix~D.
\bibitem{GL} N. Gambino and G. Lobbia, \emph{On the formal theory of
pseudomonads and pseudodistributive laws},
\href{https://arxiv.org/pdf/0907.1359v3}{arXiv:0907.1359v3},
Def.~4.1 and App.~A.
\bibitem{Lap} M. L. Laplaza,
\href{https://www.sas.rochester.edu/mth/sites/doug-ravenel/otherpapers/laplaza.pdf}
{\emph{Coherence for distributivity}}, LNM 281 (1972), 29--72, Prop.~10.
\bibitem{U} \href{https://github.com/eternal-november/eos-research-programmes/blob/f458186590df968a516ba73a8eb76fcf044544b0/projects/contextual-composition/foundations/relative-construction/delivery/One_relative_construction_for_context_content_and_effects.tex}{\emph{One relative construction for context, content and effects}},
project source at \texttt{f458186590df}.
\bibitem{O1} \href{https://github.com/eternal-november/eos-research-programmes/blob/f458186590df968a516ba73a8eb76fcf044544b0/projects/contextual-composition/interaction/independent-criterion/General_Independent_Interaction_Criterion.tex}{\emph{General independent interaction criterion}},
same pinned revision; original detailed Set proof.
\bibitem{C} \href{https://github.com/eternal-november/eos-research-programmes/blob/9e19523a6d698bbafd90987a4a49f6f3d5a27287/projects/contextual-composition/interaction/independent-criterion/Relative_Corners_and_Bispans_Comparison.tex}
{\emph{Which parts of an interactad are corners and bispans?}},
concrete comparison note, \texttt{9e19523a6d69}, 20 November 2026.
\end{thebibliography}
\end{document}
